\documentclass[11pt]{article}

\usepackage[a4paper,margin=1in]{geometry}
\usepackage{amsmath,amssymb,amsthm,mathtools}
\usepackage[T1]{fontenc}
\usepackage{lmodern}
\usepackage{microtype}
\usepackage{enumitem}
\newtheorem{theorem}{Theorem}
\newtheorem{lemma}{Lemma}
\newtheorem{proposition}{Proposition}
\newtheorem{corollary}{Corollary}

\DeclareMathOperator{\Sol}{Sol}

\title{A Self-Contained Resolution of the Integral Equation \(y(x^3-z^2)=x-1\)}
\author{}
\date{}

\begin{document}
\maketitle

\begin{abstract}
We give a completely elementary proof that the Diophantine equation
\[
    y(x^3-z^2)=x-1
\]
has infinitely many integer solutions.  In fact, we prove a precise description of the fibre above
\(x=1\), isolate the exceptional case \(x^3=z^2\), and record the exact divisibility criterion governing every remaining solution.  The infinite families obtained include infinitely many solutions with all three coordinates non-zero.
\end{abstract}

\section{Statement of the problem}

We study integer triples
\[
    (x,y,z)\in \mathbb Z^3
\]
satisfying
\begin{equation}\label{eq:main}
    y(x^3-z^2)=x-1.
\end{equation}
The goal is not to solve a finite search problem, but to decide whether there are infinitely many integer solutions.  The answer is affirmative, and the proof is short but complete.

\section{Elementary preliminary facts}

We use only the following elementary properties of the ring \(\mathbb Z\).

\begin{lemma}[Zero-product principle in \(\mathbb Z\)]\label{lem:zero-product}
If \(a,b\in \mathbb Z\) and \(ab=0\), then \(a=0\) or \(b=0\).
\end{lemma}

\begin{proof}
Suppose \(a\ne 0\).  If also \(b\ne 0\), then \(|a|\ge 1\) and \(|b|\ge 1\), so
\[
    |ab|=|a|\,|b|\ge 1,
\]
which contradicts \(ab=0\).  Hence, if \(ab=0\), at least one of \(a,b\) must be zero.
\end{proof}

\begin{lemma}[Integer square roots of \(1\)]\label{lem:square-one}
If \(z\in \mathbb Z\) and \(z^2=1\), then \(z=1\) or \(z=-1\).
\end{lemma}

\begin{proof}
The equality \(z^2=1\) is equivalent to
\[
    (z-1)(z+1)=0.
\]
By Lemma~\ref{lem:zero-product}, either \(z-1=0\) or \(z+1=0\).  Thus \(z=1\) or \(z=-1\).
\end{proof}

\section{The fibre above \(x=1\)}

The key observation is that the right-hand side of \eqref{eq:main} vanishes when \(x=1\).  This immediately produces infinite families, but we first state the exact fibre.

\begin{proposition}[Exact solution set with \(x=1\)]\label{prop:fibre-x-one}
The integer solutions of \eqref{eq:main} with \(x=1\) are precisely the triples belonging to
\[
    \{(1,0,t):t\in\mathbb Z\}
    \;\cup\;
    \{(1,t,1):t\in\mathbb Z\}
    \;\cup\;
    \{(1,t,-1):t\in\mathbb Z\}.
\]
\end{proposition}

\begin{proof}
Substituting \(x=1\) into \eqref{eq:main} gives
\[
    y(1^3-z^2)=1-1,
\]
so
\[
    y(1-z^2)=0.
\]
By Lemma~\ref{lem:zero-product}, either \(y=0\) or \(1-z^2=0\).

If \(y=0\), then \(z\) is arbitrary, and we obtain exactly the triples
\[
    (1,0,t),\qquad t\in\mathbb Z.
\]

If \(1-z^2=0\), then \(z^2=1\), and Lemma~\ref{lem:square-one} gives \(z=1\) or \(z=-1\).  In either case \(y\) is arbitrary.  Hence we obtain exactly the triples
\[
    (1,t,1),\qquad t\in\mathbb Z,
\]
and
\[
    (1,t,-1),\qquad t\in\mathbb Z.
\]

Conversely, every triple in these three displayed families satisfies \(x=1\) and either \(y=0\) or \(z^2=1\).  Therefore
\[
    y(1-z^2)=0=1-1,
\]
so each such triple satisfies \eqref{eq:main}.  This proves the claimed exact description.
\end{proof}

\section{The exceptional case \(x^3=z^2\)}

Although the proof of infinitude already follows from Proposition~\ref{prop:fibre-x-one}, it is useful to separate the only case in which the factor \(x^3-z^2\) vanishes.

\begin{lemma}[Vanishing of the cubic-square factor]\label{lem:vanishing-factor}
Let \((x,y,z)\in\mathbb Z^3\) satisfy \eqref{eq:main}.  If
\[
    x^3-z^2=0,
\]
then
\[
    x=1\quad\text{and}\quad z\in\{1,-1\}.
\]
Conversely, every triple \((1,y,1)\) and every triple \((1,y,-1)\), with \(y\in\mathbb Z\), satisfies \eqref{eq:main} and has \(x^3-z^2=0\).
\end{lemma}

\begin{proof}
Assume \((x,y,z)\) satisfies \eqref{eq:main} and \(x^3-z^2=0\).  Then the left-hand side of \eqref{eq:main} is
\[
    y(x^3-z^2)=y\cdot 0=0.
\]
Therefore the right-hand side must also be zero:
\[
    x-1=0.
\]
Thus \(x=1\).  Substituting this into \(x^3-z^2=0\) gives
\[
    1-z^2=0,
\]
so \(z^2=1\).  By Lemma~\ref{lem:square-one}, \(z=1\) or \(z=-1\).

Conversely, if \((x,z)=(1,1)\) or \((x,z)=(1,-1)\), then
\[
    x^3-z^2=1-1=0
\]
and
\[
    x-1=0.
\]
Thus, for every \(y\in\mathbb Z\), both sides of \eqref{eq:main} are zero.  Hence all such triples are solutions.
\end{proof}

\section{Divisibility criterion for the remaining cases}

The equation is linear in \(y\).  Therefore, once \(x\) and \(z\) are chosen with \(x^3\ne z^2\), the value of \(y\), if it exists, is forced.

\begin{proposition}[Exact divisibility criterion away from \(x^3=z^2\)]\label{prop:divisibility}
Let \(x,z\in\mathbb Z\) satisfy
\[
    x^3-z^2\ne 0.
\]
Then there exists an integer \(y\) such that \((x,y,z)\) satisfies \eqref{eq:main} if and only if
\[
    x^3-z^2 \mid x-1.
\]
In that case the only possible value of \(y\) is
\[
    y=\frac{x-1}{x^3-z^2}.
\]
\end{proposition}

\begin{proof}
Set
\[
    D=x^3-z^2.
\]
By hypothesis, \(D\ne 0\).  Equation \eqref{eq:main} becomes
\[
    yD=x-1.
\]

First suppose that an integer \(y\) satisfying this equation exists.  Then \(x-1\) is an integer multiple of \(D\), namely
\[
    x-1=yD.
\]
Thus \(D\mid x-1\), which is exactly
\[
    x^3-z^2\mid x-1.
\]
Moreover, since \(D\ne 0\), division by \(D\) gives the forced value
\[
    y=\frac{x-1}{D}=\frac{x-1}{x^3-z^2}.
\]

Conversely, suppose \(D\mid x-1\).  By definition of divisibility in \(\mathbb Z\), there is an integer \(q\) such that
\[
    x-1=qD.
\]
Taking \(y=q\), we get
\[
    y(x^3-z^2)=qD=x-1,
\]
so \((x,y,z)\) satisfies \eqref{eq:main}.  The displayed formula for \(y\) follows again from \(D\ne 0\).
\end{proof}

\section{Infinitely many integer solutions}

We now give the promised unconditional conclusion.

\begin{theorem}[Infinitude of integer solutions]\label{thm:infinitely-many}
The Diophantine equation
\[
    y(x^3-z^2)=x-1
\]
has infinitely many integer solutions.
\end{theorem}

\begin{proof}
For each integer \(n\), define
\[
    (x_n,y_n,z_n)=(1,n,1).
\]
Then
\[
    y_n(x_n^3-z_n^2)
    =n(1^3-1^2)
    =n(1-1)
    =n\cdot 0
    =0,
\]
while
\[
    x_n-1=1-1=0.
\]
Hence
\[
    y_n(x_n^3-z_n^2)=x_n-1,
\]
so \((1,n,1)\) is an integer solution for every \(n\in\mathbb Z\).

These solutions are pairwise distinct: if \((1,n,1)=(1,m,1)\), then equality of the second coordinates gives \(n=m\).  Since \(\mathbb Z\) is infinite, the set
\[
    \{(1,n,1):n\in\mathbb Z\}
\]
is infinite.  Therefore \eqref{eq:main} has infinitely many integer solutions.
\end{proof}

\begin{corollary}[Infinitely many non-zero-coordinate solutions]\label{cor:nonzero}
There are infinitely many integer solutions \((x,y,z)\) of \eqref{eq:main} with
\[
    xyz\ne 0.
\]
\end{corollary}

\begin{proof}
For every non-zero integer \(n\), the triple
\[
    (1,n,1)
\]
is a solution by the proof of Theorem~\ref{thm:infinitely-many}.  Its coordinate product is
\[
    1\cdot n\cdot 1=n\ne 0.
\]
Since there are infinitely many non-zero integers \(n\), this gives infinitely many solutions with all three coordinates non-zero.
\end{proof}

\section{Summary of the solution set information proved here}

Combining the preceding results, we have the following complete structural statement.

\begin{theorem}[Structural description]\label{thm:structural}
An integer triple \((x,y,z)\in\mathbb Z^3\) satisfies
\[
    y(x^3-z^2)=x-1
\]
if and only if exactly one of the following alternatives holds:
\begin{enumerate}[label=\textup{(\roman*)}]
    \item \(x=1\), \(z\in\{1,-1\}\), and \(y\in\mathbb Z\) is arbitrary;
    \item \(x=1\), \(z\notin\{1,-1\}\), and \(y=0\);
    \item \(x\ne 1\), \(x^3-z^2\ne 0\), \(x^3-z^2\mid x-1\), and
    \[
        y=\frac{x-1}{x^3-z^2}.
    \]
\end{enumerate}
In particular, alternatives \textup{(i)} and \textup{(ii)} already contain infinite families of integer solutions.
\end{theorem}

\begin{proof}
Let \((x,y,z)\in\mathbb Z^3\).

Assume first that \((x,y,z)\) is a solution.  If \(x=1\), Proposition~\ref{prop:fibre-x-one} says precisely that either \(z\in\{1,-1\}\), with \(y\) arbitrary, or else \(z\notin\{1,-1\}\), in which case the only remaining possibility is \(y=0\).  Thus a solution with \(x=1\) falls into exactly one of alternatives \textup{(i)} and \textup{(ii)}.

Now suppose \(x\ne 1\).  A solution cannot have \(x^3-z^2=0\), because Lemma~\ref{lem:vanishing-factor} would then force \(x=1\).  Hence \(x^3-z^2\ne 0\), and Proposition~\ref{prop:divisibility} gives exactly the divisibility condition and the displayed formula for \(y\).  Thus the solution falls into alternative \textup{(iii)}.

Conversely, every triple in alternative \textup{(i)} or \textup{(ii)} is a solution by Proposition~\ref{prop:fibre-x-one}, and every triple in alternative \textup{(iii)} is a solution by Proposition~\ref{prop:divisibility}.  The three alternatives are disjoint because alternatives \textup{(i)} and \textup{(ii)} require \(x=1\), alternative \textup{(iii)} requires \(x\ne 1\), and alternatives \textup{(i)} and \textup{(ii)} impose mutually exclusive conditions on \(z\).  This proves the stated structural description.
\end{proof}

\end{document}
